% GENERATED FILE. Edit canonical lessons, then run npm run generate:books.
% canonical-source-hash: fnv1a64:83567c697bfefd4b
% canonical-lessons: ZH-C357-rencuo, ZH-C357-qihua, ZH-C357-kouqi, ZH-C357-piping, ZH-C357-biaoyang, ZH-R357-review-352-357

\chapter{Owning Up}
\label{ch:zh-owning-up}
\hlchaptermodality{357}

\begin{chapteropening}
\textbf{By the end of this chapter:} \emph{I can say someone admits they were wrong, and talk about tone, angry words, blame and praise.}

Complete the last lesson of chapter 357: I can say someone admits they were wrong, and talk about tone, angry words, blame and praise.
\end{chapteropening}

\section[rèncuò]{rèncuò — to admit you were wrong}

\label{lesson:ZH-C357-rencuo}



\begin{quote}
Before the new one: say the Mandarin for an error, then the Mandarin for to put a mistake right.
\end{quote}

\subsection*{You'll want to know: rèncuò}
\emph{rèncuò} — \textquotedblleft{}to admit you were wrong\textquotedblright{}.

\begin{quote}
\emph{tā bù kěn rèncuò} — He won't admit he was wrong.
\end{quote}

Its characters come later; for now it is a word for the ear.

\begin{cousinweb}
\emph{rèncuò} is built from \emph{rèn} and \zh{错} 'wrong', both inherited from Middle Chinese, \emph{nyinH} and \emph{tshak} (Guangyun, in Baxter's transcription), which glosses the first 'to recognize'.
\end{cousinweb}

\begin{grammarlens}[title={What you've built}]
One more word for giving reasons.
\end{grammarlens}

\subsection*{Your turn}
\begin{itemize}
\raggedright
  \item \emph{Say it:} \emph{rèncuò}
  \item \emph{Say it:} \emph{rèncuò}, once more
  \item \emph{Say it:} \emph{tā bù kěn rèncuò} — He won't admit he was wrong.
  \item \emph{Recall:} say the Mandarin for an error, then the Mandarin for to put a mistake right, then say \emph{rèncuò} again
\end{itemize}

\begin{tcolorbox}[breakable,colback=teal!4,colframe=teal!35!black,title={Before you move on}]
What does rèncuò mean? (\textquotedblleft{}To admit you were wrong\textquotedblright{}.) Say it once more.
\end{tcolorbox}
\section[\texorpdfstring{\zh{气话} (qìhuà)}{气话 (qìhuà)}]{\zh{气话} (qìhuà) — something said in anger}

\label{lesson:ZH-C357-qihua}



\begin{quote}
Before the new one: say the Mandarin for to put a mistake right, then the Mandarin for to admit you were wrong.
\end{quote}

\subsection*{You'll want to know: \zh{气话}}
\textbf{\zh{气话}} — \emph{qìhuà} — \textquotedblleft{}something said in anger\textquotedblright{}.

\begin{quote}
\textbf{\zh{他说的是气话}} \emph{tā shuō de shì qìhuà} — He only said it because he was angry.
\end{quote}

\begin{cousinweb}
\textbf{\zh{气话}} (\emph{qìhuà}) is built from \zh{气} 'breath' and \zh{话} 'talk', both inherited from Middle Chinese, \emph{khj+jH} and \emph{hwaejH} (Guangyun, in Baxter's transcription), which glosses the first 'breath' and the second 'speech'.
\end{cousinweb}

\begin{grammarlens}[title={What you've built}]
One more word for giving reasons.
\end{grammarlens}

\subsection*{Your turn}
\begin{itemize}
\raggedright
  \item \emph{Say it:} \emph{qìhuà}
  \item \emph{Say it:} \emph{qìhuà}, once more
  \item \emph{Say it:} \emph{tā shuō de shì qìhuà} — He only said it because he was angry.
  \item \emph{Recall:} say the Mandarin for to put a mistake right, then the Mandarin for to admit you were wrong, then say \emph{qìhuà} again
\end{itemize}

\begin{tcolorbox}[breakable,colback=teal!4,colframe=teal!35!black,title={Before you move on}]
What does \zh{气话} mean? (\textquotedblleft{}Something said in anger\textquotedblright{}.) Say it once more.
\end{tcolorbox}
\section[\texorpdfstring{\zh{口气} (kǒuqì)}{口气 (kǒuqì)}]{\zh{口气} (kǒuqì) — tone of voice; way of saying something}

\label{lesson:ZH-C357-kouqi}



\begin{quote}
Before the new one: say the Mandarin for to admit you were wrong, then the Mandarin for something said in anger.
\end{quote}

\subsection*{You'll want to know: \zh{口气}}
\textbf{\zh{口气}} — \emph{kǒuqì} — \textquotedblleft{}tone of voice; way of saying something\textquotedblright{}.

\begin{quote}
\textbf{\zh{我觉得你说话的口气不太好}} \emph{wǒ juéde nǐ shuōhuà de kǒuqì bú tài hǎo} — I don't think your tone is very nice.
\end{quote}

\begin{cousinweb}
\textbf{\zh{口气}} (\emph{kǒuqì}) is built from \zh{口} 'the mouth' and \zh{气} 'breath', both inherited from Middle Chinese, \emph{khuwX} and \emph{khj+jH} (Guangyun, in Baxter's transcription), which glosses the first, quoting an older dictionary, 'what people speak and eat with', and the second 'breath'.
\end{cousinweb}

\begin{grammarlens}[title={What you've built}]
One more word for giving reasons.
\end{grammarlens}

\subsection*{Your turn}
\begin{itemize}
\raggedright
  \item \emph{Say it:} \emph{kǒuqì}
  \item \emph{Say it:} \emph{kǒuqì}, once more
  \item \emph{Say it:} \emph{wǒ juéde nǐ shuōhuà de kǒuqì bú tài hǎo} — I don't think your tone is very nice.
  \item \emph{Recall:} say the Mandarin for to admit you were wrong, then the Mandarin for something said in anger, then say \emph{kǒuqì} again
\end{itemize}

\begin{tcolorbox}[breakable,colback=teal!4,colframe=teal!35!black,title={Before you move on}]
What does \zh{口气} mean? (\textquotedblleft{}Tone of voice; way of saying something\textquotedblright{}.) Say it once more.
\end{tcolorbox}
\section[pīpíng]{pīpíng — to criticize}

\label{lesson:ZH-C357-piping}



\begin{quote}
Before the new one: say the Mandarin for something said in anger, then the Mandarin for a tone of voice.
\end{quote}

\subsection*{You'll want to know: pīpíng}
\emph{pīpíng} — \textquotedblleft{}to criticize\textquotedblright{}.

\begin{quote}
\emph{yīnwèi tā chídào le suǒyǐ lǎoshī pīpíng le tā} — He was late, so the teacher told him off.
\end{quote}

Its characters come later; for now it is a word for the ear.

\begin{cousinweb}
\emph{pīpíng} is built from \emph{pī} and \emph{píng}, both inherited from Middle Chinese, \emph{phej} and \emph{bjaeng} (Guangyun, in Baxter's transcription), which glosses the first 'to strike' and 'to push', and the second 'to weigh up'.
\end{cousinweb}

\begin{grammarlens}[title={What you've built}]
One more word for giving reasons.
\end{grammarlens}

\subsection*{Your turn}
\begin{itemize}
\raggedright
  \item \emph{Say it:} \emph{pīpíng}
  \item \emph{Say it:} \emph{pīpíng}, once more
  \item \emph{Say it:} \emph{yīnwèi tā chídào le suǒyǐ lǎoshī pīpíng le tā} — He was late, so the teacher told him off.
  \item \emph{Recall:} say the Mandarin for something said in anger, then the Mandarin for a tone of voice, then say \emph{pīpíng} again
\end{itemize}

\begin{tcolorbox}[breakable,colback=teal!4,colframe=teal!35!black,title={Before you move on}]
What does pīpíng mean? (\textquotedblleft{}To criticize\textquotedblright{}.) Say it once more.
\end{tcolorbox}
\section[biǎoyáng]{biǎoyáng — to praise (in front of others)}

\label{lesson:ZH-C357-biaoyang}



\begin{quote}
Before the new one: say the Mandarin for a tone of voice, then the Mandarin for to criticize.
\end{quote}

\subsection*{You'll want to know: biǎoyáng}
\emph{biǎoyáng} — \textquotedblleft{}to praise (in front of others)\textquotedblright{}.

\begin{quote}
\emph{wǒ juéde nǐ yīnggāi biǎoyáng tā} — I think you should praise him.
\end{quote}

Its characters come later; for now it is a word for the ear.

\begin{cousinweb}
\emph{biǎoyáng} is built from \emph{biǎo} and \emph{yáng}, both inherited from Middle Chinese, \emph{pjewX} and \emph{yang} (Guangyun, in Baxter's transcription), which glosses the first 'to make clear' and the second 'to raise' and 'to speak of'.
\end{cousinweb}

\begin{grammarlens}[title={What you've built}]
One more word for giving reasons.
\end{grammarlens}

\subsection*{Your turn}
\begin{itemize}
\raggedright
  \item \emph{Say it:} \emph{biǎoyáng}
  \item \emph{Say it:} \emph{biǎoyáng}, once more
  \item \emph{Say it:} \emph{wǒ juéde nǐ yīnggāi biǎoyáng tā} — I think you should praise him.
  \item \emph{Recall:} say the Mandarin for a tone of voice, then the Mandarin for to criticize, then say \emph{biǎoyáng} again
\end{itemize}

\begin{tcolorbox}[breakable,colback=teal!4,colframe=teal!35!black,title={Before you move on}]
What does biǎoyáng mean? (\textquotedblleft{}To praise (in front of others)\textquotedblright{}.) Say it once more.
\end{tcolorbox}
\section[(dialogue)]{Looking back: from Same Here to Owning Up}

\label{lesson:ZH-R357-review-352-357}



\begin{quote}
Say the Mandarin for the same feeling and for sure.
\end{quote}

\subsection*{Your turn}
\begin{itemize}
\raggedright
  \item \emph{Say it:} the same feeling, for sure, to give your approval, what one really meant, to make your view known
  \item \emph{Say it:} to talk someone round, the sense something makes, not necessarily, to argue back, to share someone's view
  \item \emph{Say it:} to put a question, a topic of conversation, a question on a test, what something is about, to have a serious talk
  \item \emph{Say it:} to decline politely, a kind thought, to be willing, to deny, to someone's face
  \item \emph{Say it:} a hard lesson learned, to miss a chance, to slip up, an error, to put a mistake right
  \item \emph{Say it:} to admit you were wrong, something said in anger, a tone of voice, to criticize, to praise
  \item \emph{Read it:} \zh{道} \zh{题} \zh{错} — each one aloud, then a word that holds it
\end{itemize}

\begin{tcolorbox}[breakable,colback=teal!4,colframe=teal!35!black,title={Before you move on}]
The same feeling? (\textbf{\zh{同感}}, \emph{tónggǎn}.) For sure? (\emph{kěndìng}.) To give your approval? (\emph{rènkě}.) What one really meant? (\textbf{\zh{本意}}, \emph{běnyì}.) To make your view known? (\emph{biǎoshì}.) To talk someone round? (\emph{shuōfú}.) The sense something makes? (\textbf{\zh{道理}}, \emph{dàoli}.) Not necessarily? (\textbf{\zh{不见得}}, \emph{bújiànde}.) To argue back? (\emph{fǎnbó}.) To share someone's view? (\emph{rèntóng}.) To put a question? (\emph{tíwèn}.) A topic of conversation? (\textbf{\zh{话题}}, \emph{huàtí}.) A question on a test? (\emph{tímù}.) What something is about? (\emph{nèiróng}.) To have a serious talk? (\emph{tánhuà}.) To decline politely? (\emph{xièjué}.) A kind thought? (\textbf{\zh{好意}}, \emph{hǎoyì}.) To be willing? (\emph{kěn}.) To deny? (\emph{fǒurèn}.) To someone's face? (\textbf{\zh{当面}}, \emph{dāngmiàn}.) A hard lesson learned? (\emph{jiàoxun}.) To miss a chance? (\textbf{\zh{错过}}, \emph{cuòguò}.) To slip up? (\textbf{\zh{出错}}, \emph{chūcuò}.) An error? (\emph{cuòwù}.) To put a mistake right? (\emph{gǎizhèng}.) To admit you were wrong? (\emph{rèncuò}.) Something said in anger? (\textbf{\zh{气话}}, \emph{qìhuà}.) A tone of voice? (\textbf{\zh{口气}}, \emph{kǒuqì}.) To criticize? (\emph{pīpíng}.) To praise? (\emph{biǎoyáng}.)
\end{tcolorbox}
